%! Characteristics of Polynomial Functions I — Unit 3, Lesson 3.4 — Slides (Beamer)
% UNITS 3+ PACKET-FIRST shape (LESSON_SHAPE.md §0a, 2026-11-14): the deck MIRRORS the Guided Notes
% frame for frame, so the class and the page never disagree. 17 frames (20 is the ceiling):
%   title → warm-up → targets → 3 × { definition → example → ladder ★/★★/★★★ → Show me | Find
%   the mistake } → Closing Practice → homework (DeltaMath online)
% Unit 3 exception: three merged skills, climbing — degree & leading coefficient → end behavior
% → symmetry. Target misconception (the crux, Round 3's Find the mistake, redacc): reading the
% degree and leading coefficient off the FIRST-WRITTEN term instead of the leading term
% (2x^2 - x^5 has degree 5, leading coefficient -1). Closing Practice 2 is its twin.
% Answers only inside \reveal{} — the handout is a Beamer handout-mode compile and drops them.
% Every graph is pre-drawn (TikZ). The deck is a work product, not part of the packet.
\documentclass[aspectratio=169, 11pt]{beamer}
\usepackage{algebra2-beamer}   % \forestheader{}, \sectionlabel[color]{}, \reveal[n]{}
\setbeamercolor{block title}{bg=forest, fg=white}
\setbeamercolor{block body}{bg=forestbg, fg=charcoal}
% Ladder rungs and Show-me options.
\newcommand{\rung}[1]{{\color{goldacc}\ifcase#1\or$\star$\or$\star\star$\or$\star\star\star$\fi}}
\newcommand{\opt}[2]{%
  \begin{tikzpicture}[baseline=-0.6ex]
    \node[circle, fill=forest, text=white, font=\bfseries\small, inner sep=1.6pt,
          minimum size=1.35em]{#1};
  \end{tikzpicture}\;#2}
% Pre-drawn graph grid: \lgrid{xmin}{xmax}{ymin}{ymax}
\newcommand{\lgrid}[4]{%
  \draw[step=1, linegray!40, very thin] (#1,#3) grid (#2,#4);
  \draw[->, semithick, charcoal] (#1,0)--(#2,0) node[below left, font=\tiny]{$x$};
  \draw[->, semithick, charcoal] (0,#3)--(0,#4) node[below left, font=\tiny]{$y$};}
% End-behavior thumbnail: \ebthumb{expression in \x}{caption}
\newcommand{\ebthumb}[2]{%
  \begin{tabular}[t]{@{}c@{}}\begin{tikzpicture}[scale=0.32]
    \draw[->, charcoal] (-2,0)--(2,0); \draw[->, charcoal] (0,-3)--(0,3);
    \begin{scope}\clip (-2,-3) rectangle (2,3);
      \draw[very thick, forest, domain=-1.9:1.9, samples=50, smooth, <->] plot (\x,{#1});
    \end{scope}
  \end{tikzpicture}\\[-1pt]{\scriptsize #2}\end{tabular}}

\begin{document}

% TITLE SLIDE (CourseName is not defined in beamer, so the name is literal)
{
\setbeamercolor{background canvas}{bg=forest}
\begin{frame}[plain]
  \vspace{2.8cm}\hspace{0.55cm}%
  \begin{minipage}{0.9\textwidth}
    {\color{goldacc}\bfseries\small LESSON 3.4}\\[0.5cm]
    {\color{white}\bfseries\LARGE Characteristics of Polynomial Functions I}\\[1.2cm]
    {\color{forestmid}\small Algebra 2~$\cdot$~Unit 3: Polynomial Functions}
  \end{minipage}
\end{frame}
}

% WARM-UP — the same three items as the warm-up page. Answers in one \reveal, ending on the
% question Round 1 answers.
\begin{frame}[t, plain]
  \forestheader{Warm-Up --- start now}
  \sectionlabel[goldacc]{3 minutes --- on your warm-up page, alone}
  \begin{enumerate}\setlength{\itemsep}{6pt}
    \item Use synthetic division to divide $x^3-2x^2-5x+6$ by $x-1$.
    \item Simplify $(-x)^2,\ (-x)^3,\ (-x)^4,\ (-x)^5$. \; If $f(x)=x^2+3$, find $f(-x)$.
    \item Write $g(x)=5x-3x^4+2-x^2$ in standard form. Name the highest power and the number
          in front of it.
  \end{enumerate}
  \vspace{2pt}
  \reveal{%
    \begin{block}{Check --- give yourself a point for each}
      1.\ $x^2-x-6$, remainder $0$ \qquad 2.\ $x^2,\ -x^3,\ x^4,\ -x^5$; \; $f(-x)=x^2+3=f(x)$
      \\[1pt]
      3.\ $-3x^4-x^2+5x+2$; \; power $4$, number $-3$ \\[3pt]
      \textbf{Hold on to this:} $5x$ was written first --- but which term controls the ends of
      the graph?
    \end{block}}
\end{frame}

% LEARNING TARGETS + HOW TODAY RUNS
\begin{frame}[t, plain]
  \forestheader{Today's targets}
  \sectionlabel{Three skills, each a step up}
  \begin{itemize}\setlength{\itemsep}{3pt}
    \item Find the \textbf{degree} and \textbf{leading coefficient} from \textbf{standard form}
          or \textbf{factored form} --- always from the \textbf{leading term}.
    \item Describe \textbf{end behavior} in arrow notation, from a graph or from the equation.
    \item Classify a function as \textbf{even}, \textbf{odd}, or \textbf{neither} --- from its
          graph and with $f(-x)$.
  \end{itemize}
  \vspace{2pt}
  \begin{block}{How today runs --- you are solving the whole hour}
    Three rounds in your Guided Notes: a short definition, one example, then \textbf{your turn}
    --- start at \rung1, climb to \rung2, and \rung3\ if you finish early --- then a
    \textbf{Show me} (1--4 fingers on the count of three) or a \textbf{Find the mistake}. We
    finish with the \textbf{Closing Practice} out of 4, then start the homework.
  \end{block}
\end{frame}

% ═══ ROUND 1 — DEGREE & LEADING COEFFICIENT ═══════════════════════════════════
\begin{frame}[t, plain]
  \forestheader{Definition: degree and leading coefficient}
  \sectionlabel{Round 1 --- read the leading term}
  \begin{block}{Degree and leading coefficient}
    A polynomial in \textbf{standard form} lists its terms from the highest power of $x$ down to
    the constant. The term with the highest power is the \textbf{leading term}: its exponent is
    the \textbf{degree}, and its coefficient is the \textbf{leading coefficient}.
  \end{block}
  \vspace{4pt}
  \begin{block}{In factored form}
    The degree is the \emph{sum} of the factors' degrees, and the leading coefficient is the
    \emph{product} of the factors' leading coefficients.
  \end{block}
  \vspace{4pt}
  \centering
  $\underbrace{-7x^5}_{\text{leading term}}+\,4x^2+x-9$
  \qquad degree $5$, \; leading coefficient $-7$
\end{frame}

\begin{frame}[t, plain]
  \forestheader{Example 1: degree and leading coefficient}
  \sectionlabel{Watch --- 2 minutes}
  \textbf{(a)} $f(x)=4x^2-7x^5+x-9$ \pause $\;=-7x^5+4x^2+x-9$ \quad \emph{(highest power
  first)} \pause \\[3pt]
  \hspace*{2em}degree $5$, leading coefficient $-7$ \quad \emph{(read the leading term $-7x^5$)}
  \pause
  \vspace{10pt}

  \textbf{(b)} $g(x)=-2(x-3)(x+1)^2$ \pause
  \[
  \begin{aligned}
    \text{degree} &= 1+2=3 && \text{add the factors' degrees} \\
    \text{leading coefficient} &= -2(1)(1)=-2 && \text{multiply leading coefficients}
  \end{aligned}
  \]
\end{frame}

\begin{frame}[t, plain]
  \forestheader{Your turn}
  \sectionlabel[goldacc]{7 minutes --- start at one star and climb}
  Find the degree and the leading coefficient.
  \begin{itemize}\setlength{\itemsep}{5pt}
    \item[\rung1] $f(x)=3x^4-2x^6+5$
    \item[\rung2] $g(x)=(2x-1)(x+4)^3$
    \item[\rung3] For the egg-drop challenge, a team folds an open-top box from a $20$\,cm by
          $12$\,cm card by cutting an $x$\,cm square from each corner:
          $V(x)=x(20-2x)(12-2x)$.
  \end{itemize}
  \vspace{2pt}
  \reveal{%
    \begin{block}{Check}
      \rung1\ $-2x^6+\dots$: degree $6$, l.c.\ $-2$ \qquad \rung2\ degree $1+3=4$, l.c.\
      $2\cdot1^3=2$ \\[1pt]
      \rung3\ degree $1+1+1=3$; l.c.\ $1\cdot(-2)\cdot(-2)=4$ (use $-2x$, not $20$ or $12$)
    \end{block}}
\end{frame}

\begin{frame}[t, plain]
  \forestheader{Show me: degree and leading coefficient of $h(x)=-x(3x+2)(x-5)^2$}
  \sectionlabel[goldacc]{Pick one --- 1 to 4 fingers on the count of three}
  \begin{columns}[T]
    \begin{column}{0.45\textwidth}
      \opt1{degree $3$, \; l.c.\ $-1$} \\[12pt]
      \opt3{degree $4$, \; l.c.\ $3$}
    \end{column}
    \begin{column}{0.45\textwidth}
      \opt2{degree $4$, \; l.c.\ $-3$} \\[12pt]
      \opt4{degree $2$, \; l.c.\ $-50$}
    \end{column}
  \end{columns}
  \vspace{10pt}
  \reveal{%
    \begin{block}{Answer: \textbf{2}}
      Degree $1+1+2=4$; l.c.\ $(-1)(3)(1)=-3$. \; \textbf{1} counted the factors, not their
      degrees, and stopped at the first one. \textbf{3} lost the leading minus.
      \textbf{4} took the biggest exponent and multiplied the constants
      $(-1)(2)(25)$.
    \end{block}}
\end{frame}

% ═══ ROUND 2 — END BEHAVIOR ═══════════════════════════════════════════════════
\begin{frame}[t, plain]
  \forestheader{Definition: end behavior}
  \sectionlabel{Round 2 --- the leading term decides the ends}
  \begin{block}{End behavior}
    The \textbf{end behavior} of a function is what $f(x)$ does at the far ends of its graph: as
    $x\to\infty$ (far right) and as $x\to-\infty$ (far left). For a polynomial, the leading term
    alone decides it. \textbf{Even} degree: both ends go the same way. \textbf{Odd} degree:
    opposite ways. \textbf{Positive} leading coefficient: the right end rises
    ($f(x)\to\infty$). \textbf{Negative}: the right end falls ($f(x)\to-\infty$).
  \end{block}
  \vspace{2pt}
  \centering
  \begin{tabular}{cccc}
  \ebthumb{1.6*(\x)^2-1.5}{even, $a>0$: both up} &
  \ebthumb{-1.6*(\x)^2+1.5}{even, $a<0$: both down} &
  \ebthumb{0.45*(\x)^3}{odd, $a>0$: down, then up} &
  \ebthumb{-0.45*(\x)^3}{odd, $a<0$: up, then down}
  \end{tabular}
\end{frame}

\begin{frame}[t, plain]
  \forestheader{Example 2: read it, then predict it}
  \sectionlabel{Watch --- 2 minutes}
  \begin{columns}[T]
    \begin{column}{0.28\textwidth}
      \centering
      \begin{tikzpicture}[scale=0.42]
        \lgrid{-3}{3}{-4}{4}
        \begin{scope}\clip (-3,-4) rectangle (3,4);
          \draw[very thick, forest, domain=-2.05:2.05, samples=80, smooth, <->]
            plot (\x,{-(\x)^4+3*(\x)^2+1});
        \end{scope}
        \node[font=\scriptsize, forest] at (-2.4,3.4) {$f$};
      \end{tikzpicture}
    \end{column}
    \begin{column}{0.68\textwidth}
      \textbf{(a)} Read the graph of $f$ at left. \pause \\[2pt]
      \hspace*{1em}as $x\to-\infty$, $f(x)\to-\infty$; \; as $x\to\infty$, $f(x)\to-\infty$
      \pause \\[2pt]
      \hspace*{1em}both ends fall $\Rightarrow$ \emph{even} degree, \emph{negative} leading
      coefficient \pause
      \vspace{10pt}

      \textbf{(b)} Predict: $g(x)=5x^2-2x^3+1$. \pause \\[2pt]
      \hspace*{1em}leading term $-2x^3$ \emph{(highest power, not first written)} \pause \\[2pt]
      \hspace*{1em}odd, negative: as $x\to-\infty$, $g(x)\to\infty$; \; as $x\to\infty$,
      $g(x)\to-\infty$
    \end{column}
  \end{columns}
\end{frame}

\begin{frame}[t, plain]
  \forestheader{Your turn}
  \sectionlabel[goldacc]{7 minutes --- start at one star and climb}
  \begin{columns}[T]
    \begin{column}{0.70\textwidth}
      Write the end behavior in arrow notation.
      \begin{itemize}\setlength{\itemsep}{5pt}
        \item[\rung1] $f(x)=x^5-4x^2+7$
        \item[\rung2] The graph of $p$ at right. Then: even or odd degree? Sign of the leading
              coefficient?
        \item[\rung3] $h(x)=(x+2)(3-x)(x-1)^2$
      \end{itemize}
    \end{column}
    \begin{column}{0.26\textwidth}
      \centering
      \begin{tikzpicture}[scale=0.36]
        \lgrid{-3}{3}{-4}{4}
        \begin{scope}\clip (-3,-4) rectangle (3,4);
          \draw[very thick, forest, domain=-2.75:2.75, samples=80, smooth, <->]
            plot (\x,{-0.5*(\x)^3+2*(\x)+1});
        \end{scope}
        \node[font=\scriptsize, forest] at (-2.4,3.4) {$p$};
      \end{tikzpicture}
    \end{column}
  \end{columns}
  \reveal{%
    \begin{block}{Check}
      \rung1\ as $x\to-\infty$, $f(x)\to-\infty$; as $x\to\infty$, $f(x)\to\infty$ \\[1pt]
      \rung2\ as $x\to-\infty$, $p(x)\to\infty$; as $x\to\infty$, $p(x)\to-\infty$ --- odd,
      negative \\[1pt]
      \rung3\ leading term $x\cdot(-x)\cdot x^2=-x^4$: both ends $\to-\infty$
    \end{block}}
\end{frame}

\begin{frame}[t, plain]
  \forestheader{Show me: the end behavior of $q(x)=-3x^5+x^2$}
  \sectionlabel[goldacc]{Pick one --- 1 to 4 fingers on the count of three}
  \begin{columns}[T]
    \begin{column}{0.45\textwidth}
      \opt1{both ends up} \\[12pt]
      \opt3{left end down, right end up}
    \end{column}
    \begin{column}{0.45\textwidth}
      \opt2{both ends down} \\[12pt]
      \opt4{left end up, right end down}
    \end{column}
  \end{columns}
  \vspace{10pt}
  \reveal{%
    \begin{block}{Answer: \textbf{4}}
      Leading term $-3x^5$: odd degree, so the ends go opposite ways; negative, so the right end
      falls. As $x\to-\infty$, $q(x)\to\infty$; as $x\to\infty$, $q(x)\to-\infty$.
      \textbf{1} and \textbf{2} treated the degree as even. \textbf{3} ignored the minus sign.
    \end{block}}
\end{frame}

% ═══ ROUND 3 — SYMMETRY (the crux) ════════════════════════════════════════════
\begin{frame}[t, plain]
  \forestheader{Definition: even and odd functions}
  \sectionlabel{Round 3 --- test with $f(-x)$}
  \begin{block}{Even, odd, or neither}
    A function is \textbf{even} if $f(-x)=f(x)$ for every $x$: its graph is symmetric about the
    $y$-axis. It is \textbf{odd} if $f(-x)=-f(x)$ for every $x$: its graph is symmetric about
    the origin (the same after a half-turn). A polynomial with all even powers (a constant is
    $x^0$) is even, one with all odd powers is odd, and a mix is \textbf{neither}.
  \end{block}
  \vspace{2pt}
  \centering
  \begin{tikzpicture}[scale=0.3]
    \begin{scope}
      \lgrid{-3}{3}{-3}{4}
      \draw[navy, dashed, semithick] (0,-3)--(0,4);
      \begin{scope}\clip (-3,-3) rectangle (3,4);
        \draw[very thick, forest, domain=-2.5:2.5, samples=80, smooth] plot (\x,{0.5*(\x)^4-2*(\x)^2});
      \end{scope}
      \node[font=\scriptsize, charcoal] at (0,-4.2) {even: $y$-axis mirror};
    \end{scope}
    \begin{scope}[xshift=11cm]
      \lgrid{-3}{3}{-3}{4}
      \fill[navy] (0,0) circle (5pt);
      \begin{scope}\clip (-3,-3) rectangle (3,4);
        \draw[very thick, forest, domain=-2.7:2.7, samples=80, smooth] plot (\x,{0.4*(\x)^3-1.6*(\x)});
      \end{scope}
      \node[font=\scriptsize, charcoal] at (0,-4.2) {odd: origin half-turn};
    \end{scope}
    \begin{scope}[xshift=22cm]
      \lgrid{-3}{3}{-3}{4}
      \begin{scope}\clip (-3,-3) rectangle (3,4);
        \draw[very thick, forest, domain=-2.6:2.6, samples=80, smooth] plot (\x,{0.3*(\x)^3+1.5});
      \end{scope}
      \node[font=\scriptsize, charcoal] at (0,-4.2) {neither};
    \end{scope}
  \end{tikzpicture}
\end{frame}

\begin{frame}[t, plain]
  \forestheader{Example 3: test it with $f(-x)$}
  \sectionlabel{Watch --- 2 minutes}
  \textbf{(a)} $f(x)=2x^4-x^2+3$ \pause
  \[
  \begin{aligned}
    f(-x) &= 2(-x)^4-(-x)^2+3 && \text{replace every $x$ with $-x$} \\
          &= 2x^4-x^2+3=f(x) && \text{even powers lose the sign: \textbf{even}}
  \end{aligned}
  \]
  \pause
  \textbf{(b)} $g(x)=x^3-4x$ \pause
  \[
  \begin{aligned}
    g(-x) &= (-x)^3-4(-x)=-x^3+4x && \text{odd powers keep the sign} \\
          &= -(x^3-4x)=-g(x) && \text{every term flipped: \textbf{odd}}
  \end{aligned}
  \]
\end{frame}

\begin{frame}[t, plain]
  \forestheader{Your turn}
  \sectionlabel[goldacc]{7 minutes --- start at one star and climb}
  \begin{columns}[T]
    \begin{column}{0.52\textwidth}
      Even, odd, or neither?
      \begin{itemize}\setlength{\itemsep}{6pt}
        \item[\rung1] Graphs A and B at right.
        \item[\rung2] $h(x)=x^5-2x^3+x$
        \item[\rung3] $k(x)=x^4+x^3$
      \end{itemize}
    \end{column}
    \begin{column}{0.44\textwidth}
      \centering
      \begin{tikzpicture}[scale=0.3]
        \begin{scope}
          \lgrid{-3}{3}{-3}{3}
          \begin{scope}\clip (-3,-3) rectangle (3,3);
            \draw[very thick, forest, domain=-2.6:2.6, samples=80, smooth] plot (\x,{-0.3*(\x)^3+0.9*(\x)});
          \end{scope}
          \node[font=\scriptsize\bfseries, charcoal] at (-2.4,2.4) {A};
        \end{scope}
        \begin{scope}[xshift=8cm]
          \lgrid{-3}{3}{-3}{3}
          \begin{scope}\clip (-3,-3) rectangle (3,3);
            \draw[very thick, forest, domain=-2.7:2.7, samples=80, smooth] plot (\x,{-0.25*(\x)^4+1.2*(\x)^2-1});
          \end{scope}
          \node[font=\scriptsize\bfseries, charcoal] at (-2.4,2.4) {B};
        \end{scope}
      \end{tikzpicture}
    \end{column}
  \end{columns}
  \vspace{4pt}
  \reveal{%
    \begin{block}{Check}
      \rung1\ A is odd (origin); B is even ($y$-axis) \\[1pt]
      \rung2\ $h(-x)=-x^5+2x^3-x=-h(x)$: odd \\[1pt]
      \rung3\ $k(-x)=x^4-x^3$ --- not $k(x)$, not $-k(x)=-x^4-x^3$: neither
    \end{block}}
\end{frame}

\begin{frame}[t, plain]
  \forestheader{Find the mistake}
  \sectionlabel[redacc]{The one that matters --- 3 minutes --- find the wrong step and fix it}
  \textbf{Student A} described $f(x)=2x^2-x^5$:
  \begin{enumerate}\setlength{\itemsep}{2pt}
    \item $f(-x)=2(-x)^2-(-x)^5=2x^2+x^5$. That is not $f(x)$ or $-f(x)$, so $f$ is
          \emph{neither}.
    \item The leading term is $2x^2$, so the degree is $2$ and the leading coefficient is $2$.
    \item Even degree, positive: as $x\to-\infty$, $f(x)\to\infty$ and as $x\to\infty$,
          $f(x)\to\infty$.
  \end{enumerate}
  \vspace{2pt}
  \reveal{%
    \begin{block}{The mistake}
      Step 2. The leading term is the \emph{highest power}, $-x^5$ --- not the term written
      first. Degree $5$, leading coefficient $-1$: odd and negative, so as $x\to-\infty$,
      $f(x)\to\infty$ and as $x\to\infty$, $f(x)\to-\infty$. Step 1 was right.
    \end{block}}
\end{frame}

% CLOSING PRACTICE — the four problems and bonus of the notes page; item 2 is the crux's twin.
\begin{frame}[t, plain]
  \forestheader{Closing Practice}
  \sectionlabel[goldacc]{6 minutes --- on your own --- score yourself out of 4}
  \begin{columns}[T]
    \begin{column}{0.78\textwidth}
      \small
      \begin{enumerate}\setlength{\itemsep}{1pt}
        \item Degree and leading coefficient of $p(x)=(3x+1)^2(2-x)$.
        \item End behavior of $q(x)=6x^3-2x^4+x$, in arrow notation.
        \item Even, odd, or neither: $r(x)=x^3+2$? Show the test.
        \item The graph of $s$ at right: end behavior, and even, odd, or neither?
      \end{enumerate}
      \textbf{Bonus (+1):} Spring-formal ticket profit at price $x$ is $P(x)=40x^2-x^3$. What
      happens to the profit if the price gets very high?
    \end{column}
    \begin{column}{0.18\textwidth}
      \centering
      \begin{tikzpicture}[scale=0.3]
        \lgrid{-3}{3}{-3}{4}
        \begin{scope}\clip (-3,-3) rectangle (3,4);
          \draw[very thick, forest, domain=-2.4:2.4, samples=80, smooth, <->] plot (\x,{0.5*(\x)^4-2*(\x)^2+1});
        \end{scope}
        \node[font=\scriptsize, forest] at (-2.3,3.4) {$s$};
      \end{tikzpicture}
    \end{column}
  \end{columns}
  \reveal{%
    \begin{block}{Check --- one point each}
      \small
      1.\ degree $3$, l.c.\ $9(-1)=-9$ \quad 2.\ leading term $-2x^4$: both ends $\to-\infty$
      \quad 3.\ $r(-x)=-x^3+2$: neither \\
      4.\ both ends $\to\infty$; even \quad \textbf{Bonus:} as $x\to\infty$, $P(x)\to-\infty$ ---
      a very high price loses money \\[1pt]
      \textbf{Watch out:} the leading term is the highest power, not the first term written ---
      and an odd degree does not make an odd function.
    \end{block}}
\end{frame}

% HOMEWORK — source: DeltaMath, online.
\begin{frame}[t, plain]
  \forestheader{Homework --- start it now}
  \sectionlabel[goldacc]{10 minutes in class, alone}
  \begin{block}{DeltaMath --- scored 1 to 5, due the first class after two study halls}
    Assignment: \textbf{A2: Lesson 3.4: End Behavior \& Symmetry}. Start it now on your device;
    the rest is due the first class after two study halls. I will announce the date in class and
    on TurtleNet.
  \end{block}
  \vspace{4pt}
  \begin{block}{Extra credit --- optional}
    The last two pages of your packet. Tear them off and turn them in with your homework.
  \end{block}
  \vspace{6pt}
  \textbf{Next:} Lesson 3.5 --- zeros and multiplicity, turning points, and where a polynomial
  rises and falls.
\end{frame}

\end{document}
